\section{Governance and Organisation}


\begin{frame}{Organisation and Governance}
    \mytext{ 
        \setlength{\parskip}{3ex}
        Governance is the system by which an organisation is controlled and operates. 
        
        Projects are one-time efforts: it follows that the organisation is created for the occasion. 
        
        But how are stakeholders, managers, partners etc. organised?
        
        
    }
    
\end{frame}    

\begin{frame}{Organisation and Governance}
    \mytext{ 
        \setlength{\parskip}{3ex}    
        Typically, a company is founded just for the purpose of the project (Special Purpose Vehicle or \textbf{SPV}).
        
        Various legal forms are possible. 
        
        Setting up, managing and evolving the organisation and its rules is a core task of project management. 
    
    \begin{itemize}
        \item In-house project: setting up a team 
        \item Otherwise: a company / association
    \end{itemize}
    
    }
    
\end{frame}    

\begin{frame}{Organisation and Governance}
    \mytext{
    \begin{itemize}
        \item Organise Relationships
        \item Define roles and responsibilities
        \item Decision making structure
        \item Steering the project
        \item Escalation management
        \item Change management
        \item Communication Strategy
    \end{itemize}
    }
\end{frame}

\begin{frame}{Projects within a Company}
    \showmytikz[1]{../shared/tikz/tikz-projectswithinacompany}
\end{frame}

\begin{frame}{Projects within a Company}
    \showmytikz[1,2]{../shared/tikz/tikz-projectswithinacompany}
\end{frame}

\begin{frame}{Projects within a Company}
    \showmytikz[2,3]{../shared/tikz/tikz-projectswithinacompany}
\end{frame}



\begin{frame}{Project Oriented Organisation}
    \showmytikz[1]{../shared/tikz/tikz-project-oriented-organisation}
\end{frame}

\begin{frame}{Project Oriented Organisation}
    \showmytikz[2]{../shared/tikz/tikz-project-oriented-organisation}
\end{frame}

\begin{frame}{Project Oriented Organisation}
    \showmytikz[1,2,3,4,5]{../shared/tikz/tikz-project-oriented-organisation}
\end{frame}




\begin{frame}{Hierarchy}
    \showmytikz{../shared/tikz/tikz-PM-smartd-hierarchy}
    \mycomment{
    

    \vfill
    The reality is more complicated!
    }
\end{frame}

\begin{frame}{Core Team Functions}
    \mytext{
        \begin{itemize}
            \item Lead(er), Manager (Product Owner?)
            \item Assistant
            \item (financial) Administration
            \item IT (support)
            \item Legal
            \item Subject Matter Experts 
            \item Partners
            \item ...
            \item Cousellor (Scrum master?)
        \end{itemize}
        
        Different ways to organise team work,
        
        different philosophies...
        
        (to be continued)
    }
\end{frame}

\begin{frame}{Governance / Rules}
    \mytext{
    
    \begin{itemize}
        \item every team should have rules
        \item typically not democratic
        \item consensus, decisions
        \item degree of rules may vary
        \item structure and clarity
        \item (psychological) safety
    \end{itemize}
    
    \medskip
    (to be continued)
    
    \medskip
    But: the team itself is also not alone!
    }
    
\end{frame}


\begin{frame}{Constitution}
    \showmytikz{../shared/tikz/tikz-constitution.tex}
\end{frame}

\begin{frame}{Project Constitution (Organisation Chart)}
    \showmytikz[caption={Project Constitution (Governance)},label=fig:projectconstitution-governance]{../shared/tikz/tikz-projectconstitution}
\end{frame}

\begin{frame}{From Project to Product: `Port'}
    \showmytikz[caption={From Project to Product: `Port' (Governance)},label=fig:port-governance]{../shared/tikz/tikz-port.tex}
\end{frame}


% \begin{frame}{Project Governance Document (PGD)}
%     \begin{itemize}
%         \item A PGD is also the `constitution' of the project
%         \item you must identify / find stakeholders
%         \item you have to create institutions
%         \item define obligations / duties
%     \end{itemize}
% \end{frame}

% \begin{frame}{From Knowledge to Action}

%     \begin{quote}
%         \flushright Der Mensch kann viel, wenn er nur muss\\
%         \flushright---unknown
%     \end{quote}

%     \begin{itemize}
%         \item to have the right to
%         \item to know how to
%         \item to have reason to
%         \item to want to
%         \item to have what it takes to
%         \item to have the duty to
%     \end{itemize}
% \end{frame}

% \begin{frame}{Obligations / Duties}
%     \begin{itemize}
%         \item bind yourself to your word
%         \item bind yourself to your supporters
%         \item bind yourself to your goals
%         \item set yourself s.m.a.r.t. goals
%         \item promise stakeholders s.m.a.r.t. goals
%         \item set team members and contractors s.m.a.r.t. goals
%         \item report and control
%     \end{itemize}
% \end{frame}


% \begin{frame}{Idea, initialisation}
%     \begin{itemize}
%         \item Why do it?
%         \item What to do to achieve what?
%         \item Who wants it?
%         \item Who pays for it?
%         \item Who might oppose it?
%         \item What will we need?
%         \item (What are the opportunities?)
%         \item (What are the risks?)
%         \item (What are the drawbacks?)
%         \item Whom do we need to convince?
%         \item To whom do we have to communicate?
%     \end{itemize}
% \end{frame}

% \begin{frame}{What Counts}
%     \begin{itemize}
%         \item Relevant for the Project Owner / Sponsor:
%         \begin{itemize}
%             \item opportunities and risks of the project
%             \item business case, business strategy
%             \item \textbf{why}
%         \end{itemize}
%         \item For the Project Manager:
%         \begin{itemize}
%             \item opportunities and risks \textbf{of tasks / sub-projects}
%             \item stakeholder satisfaction
%             \item formulate s.m.a.r.t.~goals
%             \item how
%             \item when
%             \item \textbf{who}
%             \item \textbf{with what}
%         \end{itemize}
%     \end{itemize}
% \end{frame}



% \begin{frame}{The Double Perspective of the Project Manager}
%     \begin{tikzpicture}
%         \draw  (0,0) node[below right] {Project Owner, Project Sponsors, Stake Holders} rectangle (10,-1);
        
%         \draw[->] (1,-1) -- node[right] {is entitled to $P$ and wants $P$} (1,-3);
%         \draw[<-] (6,-1) -- node[right] {is obliged to deliver $P$} (6,-3);
        
%         \draw  (0,-3) node[below right] {Project Manager w/ Steering Group} rectangle (10,-4);
        
%         \draw[->] (1,-4) -- node[right] {is entitled to $p$ and wants $p$} (1,-6);
%         \draw[<-] (6,-4) -- node[right] {is obliged to deliver $p$} (6,-6);
        
%         \draw  (0,-6) node[below right] {Project Management Office (PMO), contractors} rectangle (10,-7);
%     \end{tikzpicture}
% \end{frame}

\begin{frame}{An Example: Public Sector Project}
    \showmytikz{../shared/tikz/tikz-publicsectorproject}
\end{frame}

\begin{frame}{Example: Establishing an Organisation}
    \showmytikz{../shared/tikz/tikz-association}
\end{frame}


\begin{frame}{Initiators, Members, Partners...}
\showmytikz{../shared/tikz/tikz-initiatorsmemberspartners}


\end{frame}

\begin{frame}{PMO and Steering Group}
    \showmytikz[1]{../shared/tikz/tikz-pmoandsteeringgroup}
\end{frame}

\begin{frame}{PMO and Steering Group}
    \showmytikz[1,2]{../shared/tikz/tikz-pmoandsteeringgroup}
\end{frame}

\begin{frame}{PMO and Steering Group}
    \showmytikz[1,2,3]{../shared/tikz/tikz-pmoandsteeringgroup}
\end{frame}



\begin{frame}{Governance and Work Packages}
    \showmytikz{../shared/tikz/tikz-projectgovernance}
\end{frame}

% \begin{frame}{Project Phases}
%     \resizebox{!}{0.8\textheight}{
%         \input{tikz/tikz-chain-projectphases-details}
%     }
% \end{frame}
